% ==========================================================
%  FDM TECH - REPORTE SEMANAL (Plantilla estilo FDM)
%  Basado en el formato de C:\Users\Usuario\Desktop\moi\main(1).tex
% ==========================================================

\documentclass[11pt,a4paper]{article}

% --------------------------
% PDFLATEX FRIENDLY
% --------------------------
\usepackage[utf8]{inputenc}
\usepackage[T1]{fontenc}
\usepackage{lmodern}
\usepackage{microtype}
\usepackage[spanish]{babel}

% --------------------------
% LAYOUT (includehead evita choques header/contenido)
% --------------------------
\usepackage[
  a4paper,
  margin=1.75cm,
  top=1.95cm,
  bottom=1.85cm,
  includehead
]{geometry}

\usepackage{graphicx}
\usepackage{array}
\usepackage{tabularx}
\usepackage{booktabs}
\usepackage{enumitem}
\usepackage{setspace}
\usepackage{fancyhdr}
\usepackage{lastpage}
\usepackage{caption}
\usepackage{float}
\usepackage[hidelinks]{hyperref}
\usepackage{needspace}

\setlength{\parindent}{0pt}
\setlength{\parskip}{6pt}
\onehalfspacing

% --------------------------
% COLORES (FDM STYLE - tomado del logo)
% --------------------------
\usepackage{xcolor}
\definecolor{FDMblue}{HTML}{0A3B6F}
\definecolor{FDMcyan}{HTML}{1D9DE0}
\definecolor{FDMaccent}{HTML}{EF4337}
\definecolor{FDMdark}{HTML}{0A2E56}
\definecolor{FDMgray}{HTML}{334155}
\definecolor{FDMmuted}{HTML}{64748B}
\definecolor{FDMline}{HTML}{DCE8F3}
\definecolor{FDMsoft}{HTML}{F6FAFE}
\definecolor{FDMsoft2}{HTML}{ECF4FC}

\hypersetup{
  colorlinks=true,
  linkcolor=FDMcyan!60!FDMblue,
  urlcolor=FDMcyan!60!FDMblue,
  citecolor=FDMcyan!60!FDMblue
}

% --------------------------
% CAJAS PRO (tcolorbox)
% --------------------------
\usepackage[most]{tcolorbox}
\tcbset{
  enhanced,
  breakable,
  boxrule=0.75pt,
  colframe=FDMline,
  colback=FDMsoft,
  arc=2.2mm,
  left=4mm,right=4mm,top=3mm,bottom=3mm
}

% --------------------------
% DECORACION (barra FDM + watermark)
% --------------------------
\usepackage{tikz}
\usepackage{eso-pic}

% Barra lateral FDM
\AddToShipoutPictureBG*{%
  \begin{tikzpicture}[remember picture,overlay]
    \fill[FDMblue] (current page.north west) rectangle ([xshift=7.0mm]current page.south west);
    \fill[FDMaccent] ([xshift=7.0mm]current page.north west) rectangle ([xshift=7.7mm]current page.south west);
  \end{tikzpicture}%
}

% Watermark tenue
\AddToShipoutPictureBG{%
  \begin{tikzpicture}[remember picture,overlay]
    \node[rotate=35,scale=8,text=FDMline!55] at ([xshift=1.5cm,yshift=-3.5cm]current page.center) {FDM};
  \end{tikzpicture}%
}

% --------------------------
% ICONOS checklist
% --------------------------
\usepackage{pifont}
\newcommand{\cbox}{\ding{113}}
\newcommand{\cok}{\ding{51}}
\newcommand{\cx}{\ding{55}}

% --------------------------
% HEADERS: alturas correctas
% --------------------------
\setlength{\headheight}{44pt}
\setlength{\headsep}{12pt}

% ==========================================================
% SWITCHES
% ==========================================================
\newif\ifShowEvidence     \ShowEvidencefalse
\newif\ifShowAnnex        \ShowAnnexfalse

% ==========================================================
% VARIABLES EDITABLES
% ==========================================================
% Logos (en esta carpeta)
\newcommand{\LogoFDM}{fdm_tech_logo.png} % Banner página 1
\newcommand{\LogoHeader}{fdm_tech_logo.png} % Header página 2+

% Institucion / programa
\newcommand{\OrgLineA}{FDM Tech}
\newcommand{\OrgLineB}{FDM Tech}
\newcommand{\Ciclo}{2026-I}
\newcommand{\Programa}{FDM Tech \Ciclo}

% Documento
\newcommand{\DocTitle}{REPORTE DE ACTIVIDAD SEMANAL}
\newcommand{\Mes}{Marzo}
\newcommand{\Semana}{Semana 2}
\newcommand{\RangoFechas}{Del 09/03/2026 al 15/03/2026}

% Personas
\newcommand{\Autor}{Ruben Pi\~nas Rafael}
\newcommand{\Equipo}{Electr\'onica}
\newcommand{\Coordinador}{Direcci\'on FDM Tech}

% Horas
\newcommand{\HorasRemotas}{0}
\newcommand{\HorasPresenciales}{0}
\newcommand{\HorasTotales}{0}

% Actividades
\newcommand{\ActividadPrincipal}{Dise\~no en KiCad de la cama de sacrificio para FDM Tech, con \'area de cobre y pads para fijaci\'on, preparando el proyecto para exportar Gerber y modelado en G-code}
\newcommand{\ActividadSecundaria}{No aplica}
\newcommand{\ResumenUnaLinea}{Se complet\'o el dise\~no de la cama de sacrificio en KiCad a partir del PDF de mec\'anica, dejando listo el archivo para exportar Gerber y preparar el modelado en G-code; queda pendiente confirmar medidas del PDF.}
\newcommand{\AvancePorcentaje}{100\%}
\newcommand{\RiesgoNivel}{Medio}
\newcommand{\RiesgosDetalle}{Dependencia del PDF del \'area de mec\'anica; pueden faltar o variar medidas cr\'iticas para la ubicaci\'on de perforaciones.}
\newcommand{\RecursosNecesarios}{Validaci\'on de medidas del PDF por el \'area de mec\'anica; no se requieren otros recursos para este punto.}

% Control para "Proximos pasos" (opcional)
\newif\ifShowNextSteps      \ShowNextStepsfalse
\newcommand{\ProximoPasoUno}{No aplica.}
\newcommand{\ProximoPasoDos}{No aplica.}
\newcommand{\ProximoPasoTres}{No aplica.}

% ==========================================================
% MACROS DE ESTILO
% ==========================================================
\newcolumntype{Y}{>{\raggedright\arraybackslash}X}

\newcommand{\Field}[2]{\textbf{\color{FDMdark}#1}\hspace{4pt}{\color{FDMgray}#2}}

\newcommand{\SectionTitle}[1]{%
  \vspace{1mm}
  {\Large\bfseries\color{FDMblue}#1}\par
  \vspace{1mm}
}

\newcommand{\SubTitle}[1]{%
  {\bfseries\color{FDMdark}#1}\par
}

% Evita cortes de viñetas entre páginas
\newcommand{\SafeItem}{\Needspace{3\baselineskip}\item}

% Chip para logo blanco (banner página 1)
\newcommand{\LogoChip}[1]{%
  \tcbox[
    on line,
    colback=white,
    colframe=FDMline!85,
    boxrule=0.75pt,
    arc=3.2mm,
    left=1.8mm,right=1.8mm,
    top=1.2mm,bottom=1.2mm,
    drop shadow=black!10
  ]{#1}%
}

% Chip pequeño para header (página 2+)
\newcommand{\LogoChipSmall}[1]{%
  \tcbox[
    on line,
    enhanced,
    colback=white,
    colframe=FDMline!85,
    boxrule=0.7pt,
    arc=3.2mm,
    boxsep=0.6mm,
    left=0mm,right=0mm,
    top=0mm,bottom=0mm,
    clip upper,
    clip lower
  ]{#1}%
}

% Logo header a la derecha
\newcommand{\HeaderLogoRight}{%
  \IfFileExists{\LogoHeader}{%
    \raisebox{-0.45\height}{\LogoChipSmall{\includegraphics[height=7.6mm,trim=1.2mm 1.2mm 1.2mm 1.2mm,clip]{\LogoHeader}}}%
  }{}%
}

% URL segura (robusto para _ y #)
\newcommand{\SafeURL}{\url}

% Imagen con placeholder (robusto para nombres con _ y #)
\newcommand{\SmartFigure}[3]{%
  \begingroup
  \edef\SFfile{\detokenize{#1}}%
  \begin{center}
    \IfFileExists{\SFfile}{%
      \includegraphics[width=#2]{\SFfile}\par
      \vspace{1mm}
      {\small\color{FDMmuted}#3}%
    }{%
      \fbox{\parbox{#2}{\centering
        \vspace{10mm}
        \color{FDMmuted}\small
        Placeholder: \texttt{\SFfile}\\(Sube la imagen a la carpeta del proyecto)
        \vspace{10mm}
      }}\par\vspace{1mm}
      {\small\color{FDMmuted}#3}%
    }%
  \end{center}
  \endgroup
}

% ==========================================================
% HEADER / FOOTER
% ==========================================================
\renewcommand{\headrulewidth}{0pt}
\renewcommand{\footrulewidth}{0pt}

% Estilo página 1 (sin header; solo footer)
\fancypagestyle{firstpage}{
  \fancyhf{}
  \fancyfoot[L]{\color{FDMmuted}\small \Autor}
  \fancyfoot[C]{\color{FDMmuted}\small \OrgLineA}
  \fancyfoot[R]{\color{FDMmuted}\small Pág. \thepage\ de \pageref{LastPage}}
}

% Estilo páginas 2+ (logo pequeño arriba derecha)
\fancypagestyle{mainpage}{
  \fancyhf{}
  \fancyhead[L]{\color{FDMmuted}\small \Programa}
  \fancyhead[R]{%
    \color{FDMmuted}\small \Mes\;—\;\Semana
    \hspace{8pt}\HeaderLogoRight
  }
  \fancyfoot[L]{\color{FDMmuted}\small \Autor}
  \fancyfoot[C]{\color{FDMmuted}\small \OrgLineA}
  \fancyfoot[R]{\color{FDMmuted}\small Pág. \thepage\ de \pageref{LastPage}}
}

% Default: mainpage (para que desde pag. 2 sea automático)
\pagestyle{mainpage}

% ==========================================================
\begin{document}

% ==========================================================
% BANNER SUPERIOR (página 1)
% ==========================================================
\thispagestyle{firstpage}

\begin{tcolorbox}[
  colback=FDMdark,
  colframe=FDMdark,
  boxrule=0pt,
  arc=3.0mm,
  left=4mm,right=4mm,top=3mm,bottom=3mm
]
\begin{tabularx}{\textwidth}{@{}p{0.18\textwidth}Y@{}}

\IfFileExists{\LogoFDM}{%
  \hspace*{4mm}\raisebox{-0.30\height}{\LogoChip{\includegraphics[height=9mm]{\LogoFDM}}}%
}{%
  \hspace*{4mm}\raisebox{-0.30\height}{\LogoChip{\parbox[c][9mm][c]{0.20\textwidth}{\centering\color{FDMmuted}\small LOGO\\(\LogoFDM)}}}%
}
&
\begin{minipage}[c]{\linewidth}
  {\color{white}\bfseries\Large \DocTitle}\par
  \vspace{0.6mm}
  {\color{FDMline}\large \Programa}\par
  \vspace{0.8mm}
  {\color{FDMline}\small
    \OrgLineA \;\textbullet\;
    \Mes \;\textbullet\; \Semana \;\textbullet\; \RangoFechas
  }\par
\end{minipage}

\end{tabularx}
\end{tcolorbox}

% ==========================================================
% TARJETA EJECUTIVA
% ==========================================================
\begin{tcolorbox}[title=\textbf{\color{FDMblue}Resumen de la semana}]
\begin{tabularx}{\textwidth}{@{}p{0.24\textwidth}Y p{0.22\textwidth}Y@{}}
\toprule
\Field{Autor:}{ } & \Autor &
\Field{Area:}{ } & \Equipo \\
\Field{Direcci\'on:}{ } & \Coordinador &
\Field{Periodo:}{ } & \RangoFechas \\
\midrule
\Field{Horas remotas:}{ } & \HorasRemotas &
\Field{Horas presenciales:}{ } & \HorasPresenciales \\
\Field{Horas totales:}{ } & \HorasTotales & & \\
\bottomrule
\end{tabularx}

\vspace{2mm}
\SubTitle{Resumen}
{\color{FDMgray}\ResumenUnaLinea}

\vspace{2mm}
\SubTitle{Actividades}
\begin{itemize}[leftmargin=*,itemsep=3pt]
  \SafeItem \textbf{\color{FDMdark}Actividad principal:} \ActividadPrincipal.
  \SafeItem \textbf{\color{FDMdark}Actividad secundaria:} \ActividadSecundaria.
  \SafeItem \textbf{\color{FDMdark}Necesidades para avanzar:} \RecursosNecesarios.
\end{itemize}
\end{tcolorbox}

\clearpage

\begin{tcolorbox}[title=\textbf{\color{FDMblue}Avance de hito}]
\Field{Avance:}{ }\AvancePorcentaje
\end{tcolorbox}

\begin{tcolorbox}[title=\textbf{\color{FDMblue}Riesgos y dependencias}]\textbf{Nivel:} \RiesgoNivel\par
{\color{FDMgray}\RiesgosDetalle}
\end{tcolorbox}

\begin{tcolorbox}[title=\textbf{\color{FDMblue}Recursos necesarios}]
{\color{FDMgray}\RecursosNecesarios}
\end{tcolorbox}

\ifShowNextSteps
\begin{tcolorbox}[title=\textbf{\color{FDMblue}Proximos pasos (opcional)}]
\begin{itemize}[leftmargin=*,itemsep=2pt]
  \SafeItem \ProximoPasoUno
  \SafeItem \ProximoPasoDos
  \SafeItem \ProximoPasoTres
\end{itemize}
\end{tcolorbox}
\fi

\newpage

% ==========================================================
% 1) CONTENIDO CENTRAL
% ==========================================================
\SectionTitle{1. Contenido central (Semana 2)}
\begin{tcolorbox}
\begin{itemize}[leftmargin=*,itemsep=3pt]
  \SafeItem \textbf{Sistema FDM y subsistemas:} enfoque en electr\'onica para la cama de sacrificio.
  \SafeItem \textbf{Proceso del proyecto:} dise\~no en KiCad basado en PDF de mec\'anica y preparaci\'on para exportaci\'on.
  \SafeItem \textbf{Verificacion basica:} pendiente confirmaci\'on de medidas del PDF del \'area de mec\'anica.
\end{itemize}
\end{tcolorbox}

% ==========================================================
% 2) DESARROLLO TECNICO - PARTE I
% ==========================================================
\SectionTitle{2. Parte I — Desarrollo tecnico}
\begin{tcolorbox}[colback=FDMsoft2,title=\textbf{\color{FDMblue}Resumen tecnico}]
{\color{FDMgray}
Se dise\~no la cama de sacrificio en KiCad para FDM Tech, creando el \'area de cobre
principal y los pads correspondientes a los orificios de fijaci\'on. El archivo queda listo
para exportar Gerber y alimentar el modelado en G-code, sujeto a la validaci\'on final
de medidas del PDF de mec\'anica.
}
\end{tcolorbox}

\begin{tcolorbox}[title=\textbf{\color{FDMblue}Componentes y evidencia tecnica}]
\SubTitle{(A) Componente o modulo principal}
\begin{itemize}[leftmargin=*,itemsep=3pt]
  \SafeItem Dise\~no de cama de sacrificio en KiCad (\'area de cobre y pads de fijaci\'on).
  \SafeItem Riesgo: posibles medidas faltantes en el PDF del \'area de mec\'anica.
\end{itemize}
\SmartFigure{evidencias/kicad_cama_sacrificio_layout.png}{0.50\textwidth}{Figura 1. Layout en KiCad de la cama de sacrificio.}

\SubTitle{(B) Pruebas o validaciones}
\begin{itemize}[leftmargin=*,itemsep=3pt]
  \SafeItem No se realizaron validaciones formales en este periodo.
  \SafeItem Dependencia: confirmaci\'on de medidas del PDF de mec\'anica.
\end{itemize}
\SmartFigure{evidencias/captura_kicad_pads.png}{0.50\textwidth}{Figura 2. Captura de pads y perforaciones en KiCad.}
\end{tcolorbox}

% ==========================================================
% 3) DESARROLLO TECNICO - PARTE II
% ==========================================================
\clearpage
\SectionTitle{3. Parte II — Operacion y flujo}
\begin{tcolorbox}[colback=FDMsoft2,title=\textbf{\color{FDMblue}Flujo de trabajo observado}]
{\color{FDMgray}
Flujo seguido para convertir el PDF de mec\'anica en un dise\~no electr\'onico utilizable
para exportaci\'on Gerber y modelado en G-code.
}
\end{tcolorbox}

\begin{tcolorbox}[title=\textbf{\color{FDMblue}Pasos del proceso}]
\begin{itemize}[leftmargin=*,itemsep=3pt]
  \SafeItem Paso 1: revisar el PDF de mec\'anica y definir el \'area util de la cama.
  \SafeItem Paso 2: crear el contorno y el \'area de cobre en KiCad.
  \SafeItem Paso 3: agregar pads correspondientes a los orificios de fijaci\'on.
  \SafeItem Paso 4: preparar el proyecto para exportar Gerber y modelado en G-code.
\end{itemize}
\SmartFigure{evidencias/paso_1_pdf_mecanica.png}{0.55\textwidth}{Figura 3. Base del dise\~no a partir del PDF.}
\SmartFigure{evidencias/paso_2_area_cobre.png}{0.55\textwidth}{Figura 4. \'Area de cobre y pads de fijaci\'on.}
\end{tcolorbox}

% ==========================================================
% 4) ACTIVIDADES Y METODOLOGIA (segun ejemplo)
% ==========================================================
\SectionTitle{4. Actividades y metodologia}
\begin{tcolorbox}[title=\textbf{\color{FDMblue}Actividades}]
\begin{enumerate}[leftmargin=*,itemsep=2pt]
  \SafeItem \textbf{Actividad principal:} \ActividadPrincipal.
  \SafeItem \textbf{Actividad secundaria:} \ActividadSecundaria.
  \SafeItem \textbf{Actividad complementaria 1:} Preparaci\'on de evidencias (capturas y fotos).
  \SafeItem \textbf{Actividad complementaria 2:} Revisi\'on de consistencia con el PDF de mec\'anica.
\end{enumerate}
\end{tcolorbox}

\begin{tcolorbox}[title=\textbf{\color{FDMblue}Metodologia}]
\SubTitle{Descripcion de su actividad}
\begin{itemize}[leftmargin=*,itemsep=2pt]
  \SafeItem Preparaci\'on: revisar PDF de mec\'anica y delimitar el \'area de la cama.
  \SafeItem Ejecuci\'on: modelar en KiCad el \'area de cobre y los pads de fijaci\'on.
  \SafeItem Verificaci\'on: comparar medidas con el PDF y marcar pendientes.
  \SafeItem Cierre: dejar el proyecto listo para exportar Gerber y modelar en G-code.
\end{itemize}
\end{tcolorbox}

% ==========================================================
% 5) Referencias
% ==========================================================
\SectionTitle{5. Referencias}
\begin{tcolorbox}[title=\textbf{\color{FDMblue}Fuentes consultadas}]
\begin{itemize}[leftmargin=*,itemsep=2pt]
  \SafeItem KiCad Documentation — PCB Editor (Pcbnew). \SafeURL{https://docs.kicad.org/6.0/en/pcbnew/pcbnew.html}
  \SafeItem KiCad Documentation — Plotting/Exporting Gerber and Drill files (PCB Editor PDF). \SafeURL{https://docs.kicad.org/8.0/en/pcbnew/pcbnew.pdf}
  \SafeItem KiCad Official Website. \SafeURL{https://www.kicad.org/}
\end{itemize}
\end{tcolorbox}


\end{document}








